% !TEX root = ../main/main.tex
% Section: Ensemble Evolution - From v3.0 to v6.0 Production Deployment
% Created: November 15, 2025
% Purpose: Document 4-version evolution, failures, and production maturation

\section{Ensemble Evolution: From Research to Production (v3.0 $\to$ v6.0)}
\label{sec:ensemble-evolution}

\subsection{Overview and Motivation}

Ensemble development followed a non-linear path from academic proof-of-concept (v3.0, Oct 2025) to production deployment (v6.0, Nov 2025). This section documents all 4 major versions, including 2 failed attempts (v3.0 BNN integration, v4.0 R model integration), to provide honest accounting of research dead-ends and production lessons learned.

\paragraph{Design Philosophy Evolution.}
Early versions prioritized \textbf{model diversity} (BNN + XGBoost + State-Space) without addressing practical deployment constraints (R/Python interoperability, computational efficiency, calibration). Production versions shifted to \textbf{pragmatic heterogeneity}: models must contribute unique signal \textit{and} integrate cleanly into production infrastructure.

\paragraph{Key Insight.}
\textbf{Simple models that work beat complex models that fail.} v3.0's Bayesian Neural Network achieved 58.70 MAE but never reached production due to training instability and overfitting (see Appendix~I for full investigation). v5.1's hybrid static-incremental XGBoost (38.03 MAE) deployed successfully, demonstrating that \textit{deployment reliability > theoretical elegance}.

\subsection{Version Timeline and Performance}

Table~\ref{tab:ensemble-evolution} summarizes all 4 versions across 8 critical dimensions:

\begin{table}[htbp]
\centering
\caption{Ensemble Evolution Timeline (v3.0 $\to$ v6.0)}
\label{tab:ensemble-evolution}
\scriptsize
\begin{adjustbox}{width=\linewidth}
\begin{tabular}{@{} l l c c c c c c c @{}}
\toprule
\textbf{Version} & \textbf{Date} & \textbf{Components} & \textbf{MAE} & \textbf{ROI (\%)} & \textbf{Win Rate (\%)} & \textbf{Status} & \textbf{Deploy} & \textbf{Issue} \\
\midrule
\textbf{v3.0} & Oct 2025 & SS + BNN + Hier & 58.7 & --- & --- & ABANDONED & \xmark & BNN instability \\
\textbf{v4.0} & Oct 2025 & SS + Hier + XGB + IP & --- & --- & --- & FAILED & \xmark & rpy2 errors \\
\textbf{v5.0} & Nov 8 & SS + Hier + XGB + IP & 40.0 & 5.5 & 58.5 & BASELINE & \checkmark & --- \\
\textbf{v5.1} & Nov 14 & SS + Hier + Hybrid-XGB + IP & 39.2 & 6.0 & 59.0 & PRODUCTION & \checkmark & --- \\
\textbf{v6.0} & Nov 15 & SS + Hier + XGB-v4 + Hybrid + IP & \textbf{33.0} & \textbf{7.5} & \textbf{60.5} & A/B TEST & (pending) & --- \\
\bottomrule
\end{tabular}
\end{adjustbox}
\begin{tablenotes}
\item \textbf{Components:} SS = State-Space v2, BNN = Bayesian Neural Network (rushing only), Hier = Hierarchical v2, XGB = XGBoost v2, IP = Informative Priors v2, Hybrid-XGB = 70\% static + 30\% incremental, XGB-v4 = Ridge residual enhancement.
\item \textbf{Status:} ABANDONED = research line terminated, FAILED = integration failures prevented deployment, BASELINE = approved reference configuration, PRODUCTION = live in production, A/B TEST = deployment pending validation.
\item \textbf{v6.0 projections} based on component-level backtests (2022-2024). Actual performance to be validated via A/B test (Dec 2025).
\end{tablenotes}
\end{table}

\subsection{v3.0: BNN Integration Attempt (ABANDONED)}

\paragraph{Configuration (October 2025).}
\begin{itemize}[itemsep=2pt,parsep=0pt]
    \item \textbf{Components:} State-Space v2 (40\%), BNN v3.1 Rushing (30\%), Hierarchical v2 (30\%)
    \item \textbf{Weights:} Optimized via constrained least squares on 2023 validation
    \item \textbf{Target:} Integrate deep learning uncertainty quantification with Bayesian models
\end{itemize}

\paragraph{Failure Modes.}
\begin{enumerate}[itemsep=2pt,parsep=0pt]
    \item \textbf{BNN Passing Model Never Converged:} 6 architectures attempted over 26 days (84 person-hours), all exhibiting NaN losses or severe overfitting (see Appendix~I for full investigation)
    \item \textbf{BNN Rushing Limited Scope:} Only applicable to rushing prop bets (not passing), reducing ensemble applicability
    \item \textbf{Calibration Fragility:} BNN required temperature scaling (T=0.75) to achieve 89.5\% coverage, but calibration degraded on new seasons (72\% coverage on 2024 holdout)
\end{enumerate}

\paragraph{Strategic Decision (November 9, 2025).}
After 84 hours of debugging and 6 failed architectures, BNN research line was \textbf{terminated}. Resources redirected to hybrid XGBoost and ridge residual methods, which achieved superior performance with 10x less development time.

\paragraph{Lessons Learned.}
\begin{itemize}[itemsep=2pt,parsep=0pt]
    \item \textbf{Know when to quit:} 84 hours with no convergence = research dead-end
    \item \textbf{Simpler baselines first:} XGBoost v4 (4.70 MAE) outperformed BNN (58.70 MAE) with far less complexity
    \item \textbf{Production constraints matter:} BNN training required GPU (3.2 hours), XGBoost runs on CPU (18 minutes)
\end{itemize}

\subsection{v4.0: R Model Integration Failure (FAILED)}

\paragraph{Configuration (October 2025).}
\begin{itemize}[itemsep=2pt,parsep=0pt]
    \item \textbf{Components:} State-Space v2 (30\%), Hierarchical v2 (30\%), XGBoost v2 (25\%), Informative Priors v2 (15\%)
    \item \textbf{Weights:} Equal weighting within Bayesian (60\%) vs ML (40\%) categories
    \item \textbf{Target:} Integrate R models (brms, Stan) into Python ensemble via rpy2
\end{itemize}

\paragraph{Failure Modes.}
\begin{enumerate}[itemsep=2pt,parsep=0pt]
    \item \textbf{rpy2 Version Conflicts:} R 4.3.1 incompatible with Python 3.11.13 rpy2 bindings
    \item \textbf{Prediction Interface Errors:} R models expected \texttt{data.frame} format, Python ensemble used \texttt{numpy.ndarray}, causing type mismatches
    \item \textbf{Batch Prediction Failures:} R models processed 1 observation/call (7.9x slower than batched interface)
\end{enumerate}

\paragraph{Resolution (November 2025).}
Two-pronged fix:
\begin{enumerate}[itemsep=2pt,parsep=0pt]
    \item \textbf{Batched R Interface:} Created \texttt{r\_model\_interface\_batched.py} to process 100 predictions/call (7.9x speedup)
    \item \textbf{Environment Isolation:} Docker containers for R models to prevent version conflicts
\end{enumerate}

After fixes, v4.0 components worked individually but ensemble integration still failed due to insufficient testing. Lessons applied to v5.0 design.

\subsection{v5.0: Production Baseline (November 8, 2025)}

\paragraph{Configuration.}
\begin{table}[htbp]
\centering
\caption{Ensemble v5.0 Component Weights (Baseline)}
\label{tab:ensemble-v5-weights}
\small
\begin{tabular}{@{} l c c c c @{}}
\toprule
\textbf{Component} & \textbf{Weight (\%)} & \textbf{MAE} & \textbf{Calibration} & \textbf{Rationale} \\
\midrule
State-Space v2 & 30 & 47.77 & 92.3\% & MCMC convergence, uncertainty quantification \\
Hierarchical v2 & 30 & 52.34 & 88.7\% & Partial pooling, player-team interactions \\
XGBoost v2 (static) & 25 & 39.66 & 85.1\% & Gradient boosting, feature interactions \\
Informative Priors v2 & 15 & 53.12 & 87.4\% & Expert knowledge, domain constraints \\
\baselinerule
\textbf{Ensemble v5.0} & \textbf{100} & \textbf{40.0} & \textbf{90.2\%} & \textbf{Weighted average, temperature scaled} \\
\bottomrule
\end{tabular}
\begin{tablenotes}
\item \textbf{Notes:} Weights optimized via CVX (constrained least squares) on 2023 validation. MAE = mean absolute error on 2024 holdout (148 games). Calibration = coverage of 90\% prediction intervals. Temperature scaling (T=0.5534) applied to ensemble output for perfect 90\% calibration.
\end{tablenotes}
\end{table}

\paragraph{Validation Results (2024 Holdout).}
\begin{itemize}[itemsep=2pt,parsep=0pt]
    \item \textbf{MAE:} 40.0 yards (vs 47.77 best single model)
    \item \textbf{ROI:} 5.5\% (exceeding 5-7\% target lower bound)
    \item \textbf{Win Rate:} 58.5\% (target: 59-62\%)
    \item \textbf{Sharpe Ratio:} 1.85 (target: 2.2-2.6)
\end{itemize}

\paragraph{Production Approval (November 8, 2025).}
v5.0 passed all 5 deployment gates:
\begin{enumerate}[itemsep=2pt,parsep=0pt]
    \item $\checkmark$ MAE < 45 yards (achieved 40.0)
    \item $\checkmark$ ROI > 5\% (achieved 5.5\%)
    \item $\checkmark$ Calibration > 85\% (achieved 90.2\%)
    \item $\checkmark$ All components converged (Rhat < 1.01, ESS > 400)
    \item $\checkmark$ Temporal validation passed (3-season backtest 2022-2024)
\end{enumerate}

Status: $\checkmark$ \textbf{APPROVED BASELINE} (Nov 8, 2025).

\subsection{v5.1: Hybrid XGBoost Upgrade (November 14, 2025)}

\paragraph{Motivation.}
Incremental XGBoost investigation (Nov 10-12, 2025) revealed:
\begin{itemize}[itemsep=2pt,parsep=0pt]
    \item Pure incremental: -13.7\% underperformance vs static XGBoost
    \item Hybrid 70/30: +4.2\% improvement vs static XGBoost
    \item Rationale: Static captures long-term patterns, incremental adapts to recent trends
\end{itemize}

\paragraph{Configuration Change.}
\begin{itemize}[itemsep=2pt,parsep=0pt]
    \item \textbf{v5.0:} XGBoost v2 (static) = 25\% weight
    \item \textbf{v5.1:} Hybrid XGBoost (70\% static, 30\% incremental) = 25\% weight
    \item \textbf{Other components:} Unchanged (SS 30\%, Hier 30\%, IP 15\%)
\end{itemize}

\paragraph{Performance Improvement.}
\begin{table}[htbp]
\centering
\caption{v5.0 vs v5.1 Performance Comparison (2024 Holdout)}
\label{tab:v5-v5.1-comparison}
\small
\begin{tabular}{@{} l S[table-format=2.1] S[table-format=1.1] S[table-format=2.1] S[table-format=1.2] @{}}
\toprule
\textbf{Version} & \textbf{MAE} & \textbf{ROI (\%)} & \textbf{Win Rate (\%)} & \textbf{Sharpe} \\
\midrule
v5.0 (Baseline) & 40.0 & 5.5 & 58.5 & 1.85 \\
v5.1 (Hybrid XGB) & 39.2 & 6.0 & 59.0 & 1.92 \\
\baselinerule
\textbf{$\Delta$ (improvement)} & \textbf{-0.8} & \textbf{+0.5} & \textbf{+0.5} & \textbf{+0.07} \\
\bottomrule
\end{tabular}
\begin{tablenotes}
\item \textbf{Notes:} Improvement driven by hybrid XGBoost (+4.2\% component-level, +2\% ensemble-level after blending). All metrics improved, though incremental gains reflect 25\% XGBoost weight in ensemble.
\end{tablenotes}
\end{table}

\paragraph{Production Deployment (November 14, 2025).}
v5.1 deployed to production after A/B test validation (50 games, $p < 0.01$ improvement over v5.0). Status: $\checkmark$ \textbf{PRODUCTION} (current as of Nov 15, 2025).

\subsection{v6.0: XGBoost v4 Ridge Integration (A/B Testing)}

\paragraph{Configuration (November 15, 2025).}
\begin{table}[htbp]
\centering
\caption{Ensemble v6.0 Component Weights (5-Way)}
\label{tab:ensemble-v6-weights}
\small
\begin{tabular}{@{} l c c c c @{}}
\toprule
\textbf{Component} & \textbf{Weight (\%)} & \textbf{MAE} & \textbf{Training} & \textbf{Contribution} \\
\midrule
State-Space v2 & 27 & 47.77 & R/Stan & Uncertainty quantification \\
Hierarchical v2 & 27 & 52.34 & R/brms & Partial pooling \\
\textbf{XGBoost v4 (ridge)} & \textbf{23} & \textbf{4.70} & \textbf{Python/XGBoost} & \textbf{Ensemble residuals} \\
Informative Priors v2 & 13 & 53.12 & R/brms & Domain constraints \\
Hybrid XGBoost & 10 & 38.03 & Python/XGBoost & Incremental adaptation \\
\baselinerule
\textbf{Ensemble v6.0} & \textbf{100} & \textbf{33.0} & \textbf{---} & \textbf{Projected (pending A/B)} \\
\bottomrule
\end{tabular}
\begin{tablenotes}
\item \textbf{Notes:} Weights optimized via CVX on 2023 validation. XGBoost v4 receives highest component weight (23\%) due to 88.31\% improvement over baseline. Projected MAE assumes perfect blending; actual v6.0 performance to be validated via A/B test (Dec 2025).
\end{tablenotes}
\end{table}

\paragraph{Expected Performance (Projected).}
\begin{itemize}[itemsep=2pt,parsep=0pt]
    \item \textbf{MAE:} 33.0 yards (7 yard improvement over v5.1)
    \item \textbf{ROI:} 7.5\% (+1.5 pp vs v5.1, within 5-7\% target range)
    \item \textbf{Win Rate:} 60.5\% (+1.5 pp vs v5.1, approaching 62\% target)
    \item \textbf{Sharpe Ratio:} 2.15 (+0.23 vs v5.1, approaching 2.5 target)
\end{itemize}

\paragraph{Deployment Timeline.}
\begin{enumerate}[itemsep=2pt,parsep=0pt]
    \item \textbf{Week 1 (Nov 18-24):} A/B test launch (50/50 split, 25 games)
    \item \textbf{Week 2 (Nov 25-Dec 1):} Collect 50 games, statistical validation ($p < 0.05$)
    \item \textbf{Week 3 (Dec 2-8):} Continue to 75 games, guardrail monitoring
    \item \textbf{Week 4 (Dec 9-15):} Final analysis (100 games), deployment decision
    \item \textbf{Target Go-Live:} December 15, 2025 (pending A/B success)
\end{enumerate}

\paragraph{Risk Management.}
Guardrails for auto-stop:
\begin{itemize}[itemsep=2pt,parsep=0pt]
    \item MAE degradation > 10\% vs v5.1 (threshold: 43.1 yards)
    \item ROI decline > 2 pp vs v5.1 (threshold: 4.0\%)
    \item Calibration coverage < 80\% (target: 90\%)
\end{itemize}

\subsection{Performance Progression Across Versions}

Figure~\ref{fig:ensemble-mae-progression} shows MAE reduction across all versions:

\begin{figure}[htbp]
\centering
\begin{tikzpicture}
\begin{axis}[
    width=0.85\linewidth,
    height=6cm,
    xlabel={Ensemble Version},
    ylabel={MAE (passing yards)},
    ymin=0, ymax=65,
    xtick={1,2,3,4,5},
    xticklabels={v3.0, v4.0, v5.0, v5.1, v6.0},
    legend pos=north east,
    grid=major,
    ymajorgrids=true,
    xmajorgrids=false,
    legend style={font=\footnotesize},
    ylabel style={font=\small},
    xlabel style={font=\small},
    tick label style={font=\scriptsize}
]

% MAE progression (lower is better)
\addplot[
    color=accentblue,
    mark=*,
    line width=1.5pt
] coordinates {
    (1,58.7)  % v3.0 (BNN failed)
    (2,50)    % v4.0 (estimated, R integration failed)
    (3,40.0)  % v5.0 (baseline)
    (4,39.2)  % v5.1 (hybrid XGB)
    (5,33.0)  % v6.0 (projected)
};
\addlegendentry{Ensemble MAE}

% Best single model baseline (State-Space v2 = 47.77)
\addplot[
    color=red,
    dashed,
    line width=1pt
] coordinates {
    (1,47.77)
    (5,47.77)
};
\addlegendentry{Best Single Model (SS v2)}

% Target threshold (MAE < 35)
\addplot[
    color=green!60!black,
    dotted,
    line width=1pt
] coordinates {
    (1,35)
    (5,35)
};
\addlegendentry{Target (MAE < 35)}

% Annotations
\node[font=\scriptsize,accentblue] at (axis cs:1,60.5) {BNN failed};
\node[font=\scriptsize,accentblue] at (axis cs:2,52) {R integration};
\node[font=\scriptsize,accentblue] at (axis cs:3,42) {$\checkmark$ Baseline};
\node[font=\scriptsize,accentblue] at (axis cs:4,41.2) {$\checkmark$ Production};
\node[font=\scriptsize,accentblue] at (axis cs:5,30.5) {$\checkmark$ Projected};

\end{axis}
\end{tikzpicture}
\caption[Ensemble MAE progression across versions]{MAE progression across ensemble versions (v3.0 $\to$ v6.0). Lower is better. v3.0 and v4.0 failed to deploy due to BNN instability and R integration errors. v5.0 established production baseline (40.0 MAE). v5.1 improved via hybrid XGBoost (+2\%). v6.0 projected to achieve 33.0 MAE via XGBoost v4 ridge residuals (88.31\% component-level improvement). Dashed red line = best single model (State-Space v2, 47.77 MAE). Dotted green line = target threshold (35 MAE).}
\label{fig:ensemble-mae-progression}
\end{figure}

\subsection{Lessons Learned Across 4 Versions}

\paragraph{1. Production Constraints Dominate Academic Metrics.}
BNN (58.70 MAE) had superior theoretical foundation (variational inference, uncertainty quantification) but failed deployment due to:
\begin{itemize}[itemsep=2pt,parsep=0pt]
    \item GPU requirement (3.2 hour training vs 18 min XGBoost)
    \item Calibration fragility (89.5\% $\to$ 72\% on new seasons)
    \item Development complexity (84 hours, 6 failed architectures)
\end{itemize}

XGBoost v4 (4.70 MAE) deployed successfully with 10x less development time and superior generalization.

\paragraph{2. Incremental Improvements > Revolutionary Changes.}
v5.0 $\to$ v5.1: +2\% improvement (hybrid XGBoost), deployed successfully
v5.1 $\to$ v6.0: +15\% improvement (XGBoost v4 ridge), pending A/B test

Both upgrades preserved 77\% of existing ensemble (SS, Hier, IP unchanged), reducing deployment risk.

\paragraph{3. Multi-Season Validation Prevents Overfitting.}
XGBoost v4 validated on 3 seasons (2022-2024) with consistent performance ($\sigma = 0.40$ MAE). Single-season validation (2024 only) would miss cross-year stability.

\paragraph{4. Honest Failure Documentation Prevents Repeating Mistakes.}
Appendix~I (BNN Investigation) documents all 6 failed approaches to prevent future researchers from repeating same dead-ends. This transparency accelerates future work.

\subsection{Future Directions}

\paragraph{1. Automated Ensemble Optimization.}
Current weights hand-tuned via CVX. Future work:
\begin{itemize}[itemsep=2pt,parsep=0pt]
    \item \textbf{AutoML for ensembles:} NAS-style search over component combinations
    \item \textbf{Dynamic weighting:} Adjust weights based on recent component performance
    \item \textbf{Bayesian optimization:} Explore weight space more efficiently than grid search
\end{itemize}

\paragraph{2. Multi-Task Ensemble (Spread + Total + Props).}
Current ensemble predicts passing yards only. Extension:
\begin{itemize}[itemsep=2pt,parsep=0pt]
    \item \textbf{Multi-output models:} Predict spread, total, player props jointly
    \item \textbf{Correlation modeling:} Capture spread-total dependence for teasers
    \item \textbf{Prop bet specialization:} TD props, receptions, rushing yards
\end{itemize}

\paragraph{3. Online Ensemble Adaptation.}
Current ensemble weights static (reoptimized quarterly). Adaptive approach:
\begin{itemize}[itemsep=2pt,parsep=0pt]
    \item \textbf{Contextual bandits:} Learn which models work best for each game type
    \item \textbf{Meta-learning:} Predict ensemble weights from game metadata
    \item \textbf{Feedback loops:} Update weights based on actual game outcomes
\end{itemize}

\subsection{Chapter Summary}

Ensemble evolution from v3.0 (Oct 2025) to v6.0 (Nov 2025) demonstrates:
\begin{enumerate}[itemsep=2pt,parsep=0pt]
    \item \textbf{2 failures taught critical lessons:} BNN (84 hours wasted) and R integration (rpy2 errors) forced pragmatic design choices
    \item \textbf{v5.0 established production baseline:} 40.0 MAE, 5.5\% ROI, 90.2\% calibration (Nov 8, 2025)
    \item \textbf{v5.1 incremental improvement:} +2\% via hybrid XGBoost, deployed to production (Nov 14, 2025)
    \item \textbf{v6.0 breakthrough pending:} +15\% projected improvement via XGBoost v4 ridge residuals (Dec 2025 A/B test)
\end{enumerate}

Key takeaway: \textbf{Production reliability > theoretical elegance}. Simple models that deploy beat complex models that fail.
